% !PRINTED 467
\weberpage{467}

%% [running head, not body text -- commented out 2026-08-07] Page furniture (running
%% head, folio, signature) is excluded from this edition by standing policy.
%% \noindent \S\,127. \hfill Die Potenzen von $f(\omega)$ als Klasseninvarianten. \hfill 467

\S\,127. \textbf{\gesperrt{Die Potenzen von $f(\omega)$ als Klasseninvarianten.}}

Nach der Definition von $\gamma_2(\omega)$ ist:
\weq{1}{f(\omega)^8=\dfrac{f(\omega)^{24}-16}{\gamma_2(\omega)}.}
Ist $\omega$ die Wurzel einer Gleichung:
\weq{2}{A\omega^2+2B\omega+C=0,\quad AC-B^2=m,}
deren Diskriminante $-4m$ durch 3 nicht teilbar ist, so können wir, wenn nötig, durch Übergang zu einer äquivalenten Gleichung, $A$, $C$ als ungerade, $A$ durch 3 unteilbar und $B$ durch 3 teilbar voraussetzen. Dann aber sind nach den beiden vorigen Paragraphen $f(\omega)^{24}$ und $\gamma_2(\omega)$ Klasseninvarianten, und wir erhalten aus (1) den Satz:

3. \gesperrt{Ist $\omega$ Wurzel einer quadratischen Form von negativer, durch 3 nicht teilbarer Diskriminante, deren beide äußere Koeffizienten ungerade, deren mittlerer Koeffizient durch 6 teilbar ist, so ist}
\[
f(\omega)^8
\]
\gesperrt{eine Klasseninvariante.}

Um die Frage zu untersuchen, ob auch noch niedrigere Potenzen von $f(\omega)$ Klasseninvarianten sein können, machen wir Gebrauch von der Formel:
\weq{3}{f(\omega+2r)^6=i^{-r}f(\omega)^6,}
worin $r$ eine ganze Zahl ist; die Werte
\[
\omega_r=\omega+2r
\]
sind die Wurzeln von äquivalenten quadratischen Gleichungen
\[
A\omega_r^2+2B_r\omega_r+C_r=0,
\]
worin
\weq{4}{B_r=B-2Ar,}
und hierin läßt sich $r$ nach dem Modul 4 so bestimmen, daß bei ungeradem $m$:
\weq{5}{B_r\equiv 0,2,4,6\pmod{8}}
und bei geradem $m$:
\weq{6}{B_r\equiv 1,3,5,7\pmod{8}.}
Wenn es nun gelingt, eine Gleichung mit rationalen Koeffizienten aufzufinden, der von den vier Werten (3) entweder nur
%CONT>>
